% ===================================================================
%  TABLO No 1 — Okul. — Lise. — Sınıf.
%
%  Ceci n'est PAS une transcription : c'est une traduction, faite en
%  2026, en turc d'aujourd'hui. D'ou trois differences avec texto/io
%  et texto/fr, et trois seulement :
%
%   * pas de \nl ni de \cc — la traduction ne suit aucune ligne
%     imprimee, et les lignes de ce fichier ne sont que du confort ;
%   * pas de \begin{VUpage} — elle n'a ni feuillet ni folio, et la
%     page de lecture ne lui en affiche pas ;
%   * le gras se pose sur le TERME seul, ponctuation dehors.
%
%  Tout le reste est identique, et doit l'etre : les cles « %%K » sont
%  celles de texto/io, macro pour macro pour l'apparat, et les renvois
%  \textsuperscript{(n)} visent les MEMES objets de la planche — ce
%  sont eux qui ouvrent les gros plans.
%
%  LE TURC POSTPOSE, COMME LE HINDI, et c'est de la colonne hindie
%  qu'on est parti : verbe a la fin, cas suffixes, determinant avant
%  determine. Les seize tableaux hindis avaient deja refait leurs
%  phrases pour garder l'ordre des RENVOIS de l'ido ; il ne restait
%  qu'a changer les mots. Le renvoi appartient a son substantif et ne
%  se deplace jamais : c'est la phrase qui se refait autour de lui.
%
%  LES NOMS DE PERSONNE PRENNENT LEUR FORME TURQUE, comme dans les
%  onze autres colonnes. Or la forme turque d'un nom d'Europe est sa
%  TRANSCRIPTION : l'orthographe turque note ce qu'on prononce, et
%  c'est ainsi que l'imprimerie d'Istanbul compose depuis 1928 —
%  « Jan Jak Ruso », « Viktor Hügo ». Ioannes est donc Jan, Petrus
%  Piyer, Henrikus Anri, Stefanus Etyen, Karolus Şarl, Georgius Jorj,
%  Franciskus Fransuva. Seuls le chien Medor et l'ane Aliboron
%  gardent leur nom. Les noms latins de la note (*) ne bougent pas :
%  la note parle justement d'eux.
%
%  LE POINT SUR LE I N'EST PAS UN ORNEMENT. Le turc distingue « i »
%  de « ı » et « İ » de « I », et ce sont quatre lettres, non deux
%  casses de deux. C'est pourquoi html.py n'a pas ajoute le turc a la
%  liste commune des ordinaux : le « re.I » de Python replie « İ »
%  sur un « i » suivi d'un point combinant, et la serie cessait de se
%  reconnaitre. Les deux casses y sont ecrites a la main.
% ===================================================================

%%K t01-apar-0 apar
\VUsaut{2.0mm}
\VUcentre{12.6pt}{120}{AÇIKLAMA KİTAPÇIĞI}
\VUsaut{3.2mm}
\VUcentre{8.4pt}{160}{EKİDİR}
\VUsaut{2.6mm}
\VUtitre{15.6pt}{0}{Delmas Yardımcı Tablolarına}
\VUsaut{3.4mm}
\VUornamento{9pt}
\VUsaut{5.0mm}
\VUcentre{13.2pt}{90}{BİRİNCİ DİZİ}
\VUsaut{3.0mm}
\VUfilet{16mm}
\VUsaut{5.6mm}

%%K t01-apar-1 apar
\VUtitre{15.0pt}{60}{TABLO No 1}
\VUsaut{1.4mm}
\VUtitre{11.4pt}{0}{Okul. --- Lise. --- Sınıf.}
\VUsaut{2.2mm}
\VUfilet{20mm}
\VUsaut{6.4mm}

%%K t01-tit-1 sub
\VUpk{10.2pt}{40}{Genel açıklama.}
\VUsaut{3.0mm}

%%K t01-01-1 p
1. --- Bu tabloda bir \VUgras{lise} \VUgras{sınıfı} gösterilmiştir. Bu sınıfta
sekiz \VUgras{masa} \textsuperscript{(18)} ile sekiz \VUgras{sıra}
\textsuperscript{(32)} vardır. Her sıraya üç öğrenci oturur. \VUgras{Öğretmen}
\textsuperscript{(7)}, \VUgras{sandalyesine} \textsuperscript{(8)} oturmuş,
\VUgras{kürsüsünün} \textsuperscript{(6)} arkasındadır.

%%K t01-02-1 p
2. --- Kürsünün solunda camlı bir \VUgras{dolap} \textsuperscript{(15)} durur.
Sağında \VUgras{kara tahta} \textsuperscript{(1)} vardır; üzerinde de
\VUgras{silgisi} \textsuperscript{(2)} ile \VUgras{tebeşiri}
\textsuperscript{(3)}.

%%K t01-02-2 p
Kürsünün üstünde \VUgras{duvar levhaları} \textsuperscript{(9, 11, 12)} ile bir
\VUgras{harita} \textsuperscript{(10)} asılıdır.

%%K t01-03-1 p
3. --- Öğrencilerin hepsi \VUgras{yerinde} \textsuperscript{(17)} değildir. Jan
\textsuperscript{(4)} kara tahtanın önünde ayaktadır; elinde silgiyle
\VUgras{geometrik şekilleri} siliyor.

%%K t01-03-2 p
Piyer kürsünün yanındadır; \VUgras{yazılı ödevleri} \textsuperscript{(5)}
toplayıp öğretmene götürüyor.

%%K t01-04-1 p
4. --- Bu öğretmen \VUgras{modern diller} öğretmenidir. Öğrencilerine yabancı
dilleri \VUgras{(Almanca, İngilizce, İspanyolca, İtalyanca, Rusça} veya
\VUgras{Fransızca)} konuşmayı, okumayı ve yazmayı öğretir.

%%K t01-05-1 p
5. --- Sınıf çok geniş ve çok aydınlıktır. Dipte, hava ile ışık girsin diye,
iki \VUgras{vasistaslı} \textsuperscript{(78)} bir geniş \VUgras{pencere} ile bir
\VUgras{camlı kapı} vardır.

%%K t01-05-2 p
Bu sınıf soğuk değildir. Ortasında, onu ısıtmak için, çok iyi bir
\VUgras{soba} \textsuperscript{(46)} ile onun \VUgras{borusu}
\textsuperscript{(45)} durur.

%%K t01-05-3 p
Duvara \VUgras{askılar} \textsuperscript{(58)} çakılıdır; öğrencilerin şapkaları
ile paltoları onlara asılır.

%%K t01-tit-2 sub
\VUsaut{5.2mm}
\VUpk{10.2pt}{40}{Oyun avlusu.}
\VUsaut{3.6mm}

%%K t01-06-1 p
6. --- \VUgras{Saat} \textsuperscript{(63)} dokuzu beş geçiyor.

%%K t01-06-2 p
\VUgras{Teneffüs} \textsuperscript{(76)} yeni bitmiştir ve ders yeniden
başlıyor. Teneffüs \VUgras{avluda} \textsuperscript{(75)} olur. Birkaç öğrenciyi
oynarken görüyoruz. Bir \VUgras{yardımcı öğretmen} \textsuperscript{(74)} onları
gözetiyor.

%%K t01-07-1 p
7. --- Avluda başka kimseleri de görüyoruz : \VUgras{müfettiş}
\textsuperscript{(73)}, \VUgras{müdür} \textsuperscript{(71)} ve \VUgras{müdür
yardımcısı} \textsuperscript{(72)}.

%%K t01-07-2 p
Müfettiş okulları denetler, müdür okulu yönetir, müdür yardımcısı öğrencilerin
çalışmasını gözetir.

%%K t01-07-3 p
Dördüncü kişi \VUgras{hademedir} \textsuperscript{(77)}. Gelmeyen öğrencilerin
adlarını yazmak için koltuğunun altında bir defter taşır.

%%K t01-08-1 p
8. --- Avlunun dibinde \VUgras{jimnastik aletleri} \textsuperscript{(70)} ile
\VUgras{el işleri} \textsuperscript{(69)} atölyesi vardır.

%%K t01-08-2 p
Tablonun sağında \VUgras{müzik} ile \VUgras{resim} \VUgras{odaları} güçlükle
görünür.

%%K t01-noto-1 noto
\VUnotes{143.2mm}{%
(*) \textit{Adlar} için de bunların latince biçimleriyle yazılması öğütlenir.
Sayısız biçimden kurtulmanın tek yolu budur. Hem \textit{Giovanni,}
\textit{Jean, Johann, Jan, Hans, John, Ivan} arasından nasıl seçilsin? Adlar
için ayrı bir sözlük mü yapılsın? Latince yalın hâli alınsın ve denilsin ki :
\VUgras{Ioannes, Ioanna; Iulius, Iulia; Iakobus; Andreas; Lukas.}
(Tam dil bilgisi).%
}

%%K t01-tit-3 sub
\VUpk{10.2pt}{40}{Ayrıntılı açıklama.}
\VUsaut{2.4mm}

%%K t01-tit-4 sub
\VUpk{10.2pt}{40}{İyi öğrenciler.}
\VUsaut{3.0mm}

%%K t01-09-1 p
9. --- Kürsünün sağında, birinci sıraya oturmuş öğrenciler \VUgras{Anri} (*),
\VUgras{Etyen} ve \VUgras{Şarl'dır}.

%%K t01-09-2 p
Anri çok dikkatli ve çalışkandır; öğretmeni dinliyor. Masasında bir
\VUgras{defter} \textsuperscript{(28)}, bir \VUgras{kitap}
\textsuperscript{(29)}, bir \VUgras{sözlük} \textsuperscript{(30)} ve bir
\VUgras{kalemlik} \textsuperscript{(31)} vardır. Kitabında pek çok
\VUgras{sayfa} vardır; her bölümde birçok \VUgras{paragraf}, her
\VUgras{satırda} da pek çok \VUgras{kelime}.

%%K t01-09-4 p
Yanındaki Etyen \textsuperscript{(24)} heveslidir. Gelecek sefere verilen dersi
defterine yazıyor. \VUgras{Yazısı} güzeldir. Kitabı kapalıdır. Yalnız
\VUgras{cildi} \textsuperscript{(27)} görünür. Yanında \VUgras{pergeli}
\textsuperscript{(26)} durur.

%%K t01-09-5 p
Şarl \textsuperscript{(23)} kitabı elinde tutuyor; onu açıp dersini yeniden
tekrarlıyor. Her öğrenci gibi onun önünde de bir \VUgras{mürekkep hokkası}
\textsuperscript{(25)} vardır. Söz dinler.

%%K t01-10-1 p
10. --- Sonraki masada \VUgras{Lui, Antuan} ile \VUgras{Pol} vardır. Lui
\VUgras{dersini} \textsuperscript{(22)} anlatıyor ve öğretmenin
\VUgras{sorularına} karşılık veriyor. Antuan \textsuperscript{(21)} çok
dikkatsizdir; öğretmen ona kimi zaman ceza da verir. Pol
\textsuperscript{(20)} iyi bir arkadaştır, uysal ve yardımseverdir. Çok dikkat
eder.

%%K t01-tit-5 sub
\VUsaut{4.2mm}
\VUpk{10.2pt}{40}{Kötü öğrenciler.}
\VUsaut{3.2mm}

%%K t01-11-1 p
11. --- Anri'nin arkasında \VUgras{Jorj} \textsuperscript{(38)} oturur. Kötü bir
çocuk değildir, ama \VUgras{gevezedir}, dikkatsiz ve yerinde duramazdır.
\VUgras{Hoppalığı} ile \VUgras{söz dinlemezliği} ders sırasında
\VUgras{karışıklık} çıkarır, ve ona sık sık sert bir \VUgras{uyarı} gerekir.
\VUgras{Müsvedde defteri} \textsuperscript{(35)} kirlidir, \VUgras{kalemiyle}
\textsuperscript{(34)} onun üzerine \VUgras{mürekkep lekeleri} damlatır durur;
bu yüzden \VUgras{silgisi} \textsuperscript{(36)} her vakit gerekir;
\VUgras{çantası} \textsuperscript{(33)} yırtıktır, çünkü çok özensizdir. Modern
diller sınıfına \VUgras{kalem tutacağını} \textsuperscript{(37)} niçin getirir?

%%K t01-11-2 p
Yanındaki Edmon \textsuperscript{(39)} ondan hoşlanmaz. Doğrusu biraz
tembeldir, ama sakin ve durgundur. Gevezelik etmez; yalnız, dinlemek yerine
\VUgras{ders kitabının} \VUgras{hikâyelerini} okumayı yeğler.

%%K t01-11-3 p
Bu sıranın üçüncü öğrencisi \VUgras{Lüdovik'tir} \textsuperscript{(40)}.
Çalışkandır. Beyaz bir \VUgras{kâğıt} \VUgras{yaprağına} yazıyor. Önüne
\VUgras{kurutma kâğıdını} \textsuperscript{(41)} koymuştur.

%%K t01-12-1 p
12. --- Öğretmenin solunda, ikinci sıranın öğrencileri pek küçüktür. İlki,
küçük \VUgras{Emil}, daha iyi yazamıyor; \VUgras{taş tahtasını}
\textsuperscript{(42)}, \VUgras{taş kalemini} ve \VUgras{süngerini}
\textsuperscript{(43)} getirir.

%%K t01-12-2 p
Yanında \VUgras{Ogüst} ile \VUgras{Bernar} \textsuperscript{(44)} vardır.
Sonuncusu iyi çalışır, ama \VUgras{kılığı} pek pasaklıdır.

%%K t01-13-1 p
13. --- Arkalarında üç \VUgras{tembel} oturur. \VUgras{Alber} derslerini
öğrenmez. \VUgras{Jak} \textsuperscript{(47)} dinlemek yerine uyuyor.
\VUgras{Tembelliği} ona \VUgras{azar} ile \VUgras{ceza} bile getirir. Sonunda
\VUgras{Kristiyan} \textsuperscript{(48)} çok daha yerinde duramazdır.
\VUgras{Davranışı} kötüdür ve düzelmek için hiçbir çaba göstermez.

%%K t01-14-1 p
14. --- Yanlarındakiler ise tersine uslu çocuklardır. \VUgras{Ernest}
\textsuperscript{(51)} düzeni ve kuralı sever; her vakit \VUgras{cetvelini}
\textsuperscript{(50)} ile \VUgras{kurşun kalemini} \textsuperscript{(49)}
kullanır. \VUgras{Gündüzlü öğrencidir}.

%%K t01-14-2 p
\VUgras{Fransuva} \textsuperscript{(52)} \VUgras{yatılı öğrencidir}; üniforma
giyer, okulda yer ve okulda yatar. İyi bir \VUgras{arkadaştır}.

%%K t01-14-3 p
Moris \VUgras{çakısıyla} \textsuperscript{(56)} \VUgras{kurşun kalemini}
yontuyor; çok özenlidir.

%%K t01-14-4 p
Bütün kitaplarını getirir, \VUgras{dil bilgisini} \textsuperscript{(55)},
\VUgras{not defterini} \textsuperscript{(54)}, bir de \VUgras{yazı altlığı}
\textsuperscript{(53)}. \VUgras{Torbası} \textsuperscript{(57)} arkasında yerde
durur.

%%K t01-14-5 p
Sınıfın arkasındaki iki masaya \VUgras{iyi öğrenciler} \textsuperscript{(19)}
oturmuştur.

%%K t01-tit-6 sub
\VUsaut{4.6mm}
\VUpk{10.2pt}{40}{Resim ve müzik.}
\VUsaut{3.4mm}

%%K t01-15-1 p
15. --- \VUgras{Resim odasında} duvara güzel \VUgras{örnekler} asılıdır ve
tavandan \VUgras{abajurlu} bir \VUgras{lamba} \textsuperscript{(62)} sarkar.
\VUgras{Öğretmen} \textsuperscript{(60)} pek sabırlıdır; öğrencilerin
\VUgras{resimlerini} \textsuperscript{(61)} düzeltir ve önlerine konulan
\VUgras{heykele} \textsuperscript{(59)} bakarak resim yapmayı onlara öğretir.

%%K t01-16-1 p
16. --- \VUgras{Müzik odası} pek ilgi çekicidir. Orada \VUgras{müzik} okumayı,
\VUgras{portede} yazılı \VUgras{notaları} \textsuperscript{(67)} söylemeyi (do,
re, mi, fa, sol, la, si\,=\,C. D. E. F. G. A. B.), \VUgras{anahtarları}
tanımayı ve hoş \VUgras{ezgiler} söylemeyi öğretirler. Notalar \VUgras{nota
sehpasına} \textsuperscript{(65)} konur. Bir öğrenci \VUgras{piyano çalıyor}
\textsuperscript{(64)} ve \VUgras{müzik öğretmeni} \textsuperscript{(66)}
\VUgras{tempo tutuyor} \textsuperscript{(68)}.

%%K t01-17-1 p
17. --- \VUgras{El işleri} \textsuperscript{(69)} \VUgras{atölyesinin} bir
köşesini güçlükle görüyoruz; çocuklar orada tahta yontmayı ve demir eğelemeyi
öğrenebilirler. Bu her okulda bulunmaz.

%%K t01-tit-7 sub
\VUsaut{5.0mm}
\VUpk{10.2pt}{40}{Fen odası.}
\VUsaut{3.6mm}

%%K t01-18-1 p
18. --- Matematik öğretmeni öğrencilere \VUgras{geometri, aritmetik, hesap} ve
\VUgras{metrik sistem} okutur. Tabloda geometrinin başlıca şekillerini
görüyoruz : bir \VUgras{daire} \textsuperscript{(a)}, bir \VUgras{kare}
\textsuperscript{(b)}, bir \VUgras{üçgen} \textsuperscript{(c)}, bir
\VUgras{doğru} \textsuperscript{(d)}.

%%K t01-18-2 p
Bir de \VUgras{işlem} görüyoruz; ne \VUgras{toplamadır}, ne
\VUgras{çıkarmadır}, ne \VUgras{çarpmadır} : \VUgras{bölmedir}
\textsuperscript{(e)}.

%%K t01-19-1 p
19. --- Metrik sistem tablosunda şunları görüyoruz : bir \VUgras{metre}
\textsuperscript{(a)}, bir \VUgras{terazi} \textsuperscript{(b)} ile
\VUgras{ağırlıkları}, bir \VUgras{litre} \textsuperscript{(e)}, bir
\VUgras{dekalitre} \textsuperscript{(f)}, bir \VUgras{desilitre} ve bir
\VUgras{metreküp} \textsuperscript{(d)}.

%%K t01-19-2 p
Öğrenciler \VUgras{doğa bilimlerini} \textsuperscript{(11)} --- zooloji
\textsuperscript{(a)}, botanik \textsuperscript{(b)}, jeoloji
\textsuperscript{(c)} --- ve \VUgras{tarihi} \textsuperscript{(12)} de okurlar.

%%K t01-19-3 p
Dolapta \VUgras{fizik} \textsuperscript{(15)} ile \VUgras{kimya}
\textsuperscript{(16)} aletleri durur.

%%K t01-19-4 p
Dolabın üstünde şunlar vardır : bir \VUgras{küre} \textsuperscript{(14)}, bir
\VUgras{koni} ve bir \VUgras{yerküre} \textsuperscript{(13)}.

%%K t01-19-5 p
Tabloların üstünde, dünyanın beş parçasını gösteren bir \VUgras{harita}
asılıdır : \VUgras{Avrupa} \textsuperscript{(g)}, \VUgras{Asya}
\textsuperscript{(e)}, \VUgras{Afrika} \textsuperscript{(f)},
\VUgras{Amerika} \textsuperscript{(ab)} ve \VUgras{Okyanusya}
\textsuperscript{(h)}, bir de iki büyük okyanus, \VUgras{Büyük Okyanus}
\textsuperscript{(c)} ile \VUgras{Atlas Okyanusu} \textsuperscript{(d)}.
\VUsaut{6.0mm}
\VUfilet{22mm}
